% Chapter 1 -- Getting started (model: TP 6.0 UG Ch 1 "Learning the new
% IDE", AN-19 tutorial: the first program typed, compiled, run; numbered
% steps; margin notes; ends on a forward pointer).  The IDE is not shipped
% (ADR-0007 DEC-08, M8), so the tutorial uses the command-line compiler.
% Sources (Law 1): os/tps/tps.pas (fixed TPSIN.PAS / TPSOUT.S, the four
% bracketed regions); the Makefile gate test-tps-gen-os (TPS.COM typed as
% `tps' at the InitechDOS prompt generates TPSOUT.S on the OS; the host
% assembles and links; the .COM runs on InitechDOS); seed/rt/start_dos.asm
% (output on COM1).  The compiler output shown is replayed by test-pubs-turbo
% through tps.pas built by Free Pascal; test-tps-lex-os, -parse-os,
% -type-os and -gen-os prove TPS.COM prints the same on InitechDOS.
\Gates{test-tps-lex-os, test-tps-parse-os, test-tps-type-os, test-tps-gen-os, test-compiler-os, test-pubs-turbo}
\chapter{Getting started}
\ix{getting started}

In this chapter you write a program, compile it with Turbo Initech, and
run it.  Along the way you meet everything the compiler tells you about a
program.  By the end of the chapter you will know the whole cycle that
every Turbo Initech program goes through.

\section{The cycle}
\ix{compile cycle}
A Turbo Initech program passes through four steps:

\begin{enumerate}
\item You write the program as a text file.
\item The compiler, \file{TPS.COM}, reads the program from the file
  \file{TPSIN.PAS} and writes it as assembly language to the file
  \file{TPSOUT.S}.
\item The assembly language is assembled and linked with the Turbo
  Initech run-time library into a program file, such as \file{FIRST.COM}.
\item You run the program at the InitechDOS prompt.
\end{enumerate}

\MarginNote{TPS always reads TPSIN.PAS.}The compiler does not take a file
name: it always reads \file{TPSIN.PAS} in the current directory and always
writes \file{TPSOUT.S}.  Keep each program in a file of its own name and
copy it to \file{TPSIN.PAS} when you compile it.

\section{Writing your first program}
\ix{first program}
\begin{enumerate}
\item Start the editor.  With the InitechDOS line editor, type
  \typed{EDLIN FIRST.PAS} and press \key{Enter}.
\item Type the program exactly as shown:

\Program{first}

\item Save the file and leave the editor.  In EDLIN, the command \typed{I}
  inserts the lines you type, \key{Ctrl-C} ends the insertion, and the
  command \typed{E} saves the file and ends EDLIN (\doctitle{InitechDOS
  Reference}, EDLIN).
\end{enumerate}

\rw{program} names the program; the statements between \rw{begin} and
\rw{end} are what it does; the period ends it.  \std{writeln}{Writeln}
writes the text between the apostrophes and ends the line.

\section{Compiling the program}
\ix{compiling}
\begin{enumerate}
\item Copy the program to the compiler's input file: type
  \typed{COPY FIRST.PAS TPSIN.PAS} and press \key{Enter}.
\item Type \typed{TPS} and press \key{Enter}.
\end{enumerate}

The compiler reads the program four times, each time doing more with it,
and reports each reading on COM1 between a pair of lines.  The first part
lists the tokens of the program, the words and symbols it is made of:

\CompilerOutput{first}{LEX}

The second part traces how the compiler recognized the structure of the
program: it entered the \textit{program} rule on line 1, the
\textit{block} on line 2, and so on, and it ends with \code{TPS-PARSE-OK}:

\CompilerOutput{first}{PARSE}

The third part lists the names the program declares.  This program
declares none, so the list is empty:

\CompilerOutput{first}{TYPE}

The last part reports the assembly language written to \file{TPSOUT.S}:
587 bytes, with one label.

\CompilerOutput{first}{GEN}

\MarginNote{Look for the -OK lines.}When all four parts end with
\code{-OK}, the program has compiled.  When the compiler finds an error, it
prints a line ending in \code{-ERROR} that says where the error is and what
it is, and does not write \file{TPSOUT.S}; Chapter 4 shows you how to read
these lines.

\section{Making the program file}
\ix{assembling}
\ix{linking}
\file{TPSOUT.S} is assembled and linked with the Turbo Initech run-time
library into the program file \file{FIRST.COM}.  The
\doctitle{Programmer's Guide}, Chapter 15, describes the file and how it is
linked.
\draft{an InitechOS-resident assembler and linker are not built: TPSOUT.S
is assembled and linked on the development system (Self-Compilation Guide,
Chapter 3); milestone M8, bead initech-ls4}

\section{Running the program}
\ix{running a program}
\begin{enumerate}
\item At the InitechDOS prompt, type \typed{FIRST} and press \key{Enter}.
\end{enumerate}

The program writes to COM1:

\ProgramOutput{first}

and InitechDOS shows its prompt again.

\section{Where to go from here}
You have written, compiled and run a program.  Chapter 2 teaches the
language: numbers, decisions, loops, characters and strings, arrays,
records, routines and files, each with a short program to compile and run.
